\documentclass{article}
\usepackage{graphicx} % Required for inserting images
\usepackage{amssymb}

\title{Aparcament automàtic}
\author{David Elias, Álvaro Noya, Xavier Parellada, Guillem Sala}
\date{March 2026}

\begin{document}
\maketitle

\section{Definicions}
\begin{enumerate}
    \item \textbf{Pas:} Anomenem pas $p$ a la parella $p=(R,d)$ formada pel radi $R\in\mathbb{R}\cup[-\infty,\infty]$ i la distància $d\in\mathbb{R}$.
    \item \textbf{Maniobra:} conjunt de passos. (Segurament ordenats amb un índex?)
    \item \textbf{Cotxe:} Pensem el cotxe com quatre punts $C:=((a_1,b_1),(a_2,b_2),(a_3,b_3),(a_4,b_4))$, on $a_i,b_i\in R$. La llargada $L$ és la llargada del cotxe ve donada per $L=dist((a_1,b_1),(a_3,b_3))$ i la seva amplada és $l=dist((a_1,b_1),(a_2,b_2))$. 
    \item \textbf{Centre del cotxe:} baricentre del rectangle cotxe. 
    \item \textbf{Llargada aparcament:} Llargada aparcament $L_A$. Es considerarà un lloc "aparcable" si $L_A\leq L+x$
    \item \textbf{Distància mínima $min(L_A)$:}
    \item \textbf{Distància raonable $L_R$:}
\end{enumerate}

\section{Moviments del cotxe}
    \begin{itemize}
        \item Anar endavant
        \item Anar endarrere
        \item Moviment horitzontal

    \end{itemize}


\section{Hipòtesis simplificadores}


\end{document}
